\documentclass[12pt]{article}

\usepackage{amsmath, amssymb}
\usepackage{geometry}
\usepackage{array}
\usepackage{booktabs}
\usepackage{physics}
\usepackage{hyperref}

\geometry{margin=1in}

\title{Assignment 1}
\author{Zachary Tullier}
\date{6/8/26}

\begin{document}

\maketitle

\section*{Q1. What are radiations? Do they satisfy the properties of waves?}

\textbf{Radiation} is energy emitted and transported through space or matter. In physics, radiation can appear as:

\begin{enumerate}
    \item \textbf{Electromagnetic radiation}, such as radio waves, microwaves, infrared, visible light, ultraviolet, X-rays, and gamma rays.
    \item \textbf{Particle radiation}, such as alpha particles, beta particles, neutrons, protons, electrons, and neutrinos.
    \item \textbf{Gravitational radiation}, or gravitational waves, produced by accelerating masses.
\end{enumerate}

Radiation can behave like a \textbf{wave}, a \textbf{particle}, or both, depending on the situation. Electromagnetic radiation especially shows \textbf{wave-particle duality}.

Electromagnetic radiation satisfies the main properties of waves:

\begin{itemize}
    \item It has wavelength.
    \item It has frequency.
    \item It has speed.
    \item It can undergo reflection, refraction, diffraction, interference, and polarization.
    \item It transports energy and momentum.
\end{itemize}

In vacuum, electromagnetic radiation is described by Maxwell's equations. A plane electromagnetic wave may be written as

\begin{equation}
\mathbf{E}(\mathbf{r},t)=\mathbf{E}_0\cos(\mathbf{k}\cdot\mathbf{r}-\omega t),
\end{equation}

\begin{equation}
\mathbf{B}(\mathbf{r},t)=\mathbf{B}_0\cos(\mathbf{k}\cdot\mathbf{r}-\omega t),
\end{equation}

where $\mathbf{E}$ is the electric field, $\mathbf{B}$ is the magnetic field, $\mathbf{k}$ is the wave vector, and $\omega$ is the angular frequency.

For electromagnetic waves in vacuum,

\begin{equation}
c=\frac{1}{\sqrt{\mu_0\varepsilon_0}},
\end{equation}

where $c$ is the speed of light, $\mu_0$ is the permeability of free space, and $\varepsilon_0$ is the permittivity of free space.

Quantum mechanically, electromagnetic radiation is made of photons. Each photon has energy

\begin{equation}
E=h\nu=\hbar\omega.
\end{equation}

So radiation behaves as a wave during propagation and as particles during emission, absorption, and detection.

\section*{Q2. Types of radiations and how they differ}

Radiation can be classified in several ways.

\subsection*{1. Ionizing vs. non-ionizing radiation}

\textbf{Non-ionizing radiation} does not usually have enough energy per quantum to remove tightly bound electrons from atoms. Examples include

\[
\text{radio waves, microwaves, infrared, visible light, low-energy ultraviolet.}
\]

\textbf{Ionizing radiation} has enough energy to ionize atoms or molecules. Examples include

\[
\text{high-energy ultraviolet, X-rays, gamma rays, alpha particles, beta particles, neutrons.}
\]

The key distinction is photon or particle energy. Ionization typically requires energies of order several electronvolts or more, depending on the atom or molecule.

\subsection*{2. Electromagnetic radiation}

Electromagnetic radiation consists of oscillating electric and magnetic fields. It does not require a material medium.

Types include

\[
\text{radio} \rightarrow \text{microwave} \rightarrow \text{infrared} \rightarrow \text{visible} \rightarrow \text{ultraviolet} \rightarrow \text{X-ray} \rightarrow \gamma\text{-ray}.
\]

They differ mainly by \textbf{frequency}, \textbf{wavelength}, and \textbf{energy}.

Higher frequency means shorter wavelength and higher photon energy:

\begin{equation}
E=h\nu=\frac{hc}{\lambda}.
\end{equation}

Gamma rays and X-rays are high-energy electromagnetic radiations. Radio waves are low-energy electromagnetic radiations.

\subsection*{3. Particle radiation}

Particle radiation consists of moving particles with mass or charge.

Important examples are shown in Table~\ref{tab:radiation_types}.

\begin{table}[h!]
\centering
\begin{tabular}{0.8\textwidth}
\toprule
\textbf{Radiation type} & \textbf{Particle} & \textbf{Charge} & \textbf{Main origin} & \textbf{Penetration} \\
\midrule
Alpha radiation & $^4\mathrm{He}^{2+}$ nucleus & $+2e$ & Nuclear decay & Low \\
Beta-minus radiation & Electron & $-e$ & Weak nuclear decay & Moderate \\
Beta-plus radiation & Positron & $+e$ & Weak nuclear decay & Moderate \\
Neutron radiation & Neutron & $0$ & Fission, fusion, nuclear reactions & High \\
Proton radiation & Proton & $+e$ & Cosmic rays, accelerators & Moderate to high \\
Neutrino radiation & Neutrino & $0$ & Weak interactions, stars, reactors & Extremely high \\
\bottomrule
\end{tabular}
\caption{Examples of particle radiation.}
\label{tab:radiation_types}
\end{table}

For massive particles, the energy is not simply $E=h\nu$. Instead, the relativistic energy is

\begin{equation}
E^2=p^2c^2+m^2c^4.
\end{equation}

For photons, $m=0$, so

\begin{equation}
E=pc.
\end{equation}

For massive particles, their de Broglie wavelength is

\begin{equation}
\lambda=\frac{h}{p}.
\end{equation}

Thus even particle radiation can have wave-like behavior.

\subsection*{4. Gravitational radiation}

Gravitational radiation consists of propagating disturbances in spacetime curvature. These are gravitational waves.

In linearized general relativity, the metric is written as

\begin{equation}
g_{\mu\nu}=\eta_{\mu\nu}+h_{\mu\nu},
\end{equation}

where $\eta_{\mu\nu}$ is the flat spacetime metric and $h_{\mu\nu}$ is a small perturbation. In vacuum, gravitational waves satisfy a wave equation:

\begin{equation}
\Box h_{\mu\nu}=0.
\end{equation}

where the d’Alembertian operator, or wave operator, is written as

\begin{equation}
\Box \equiv \frac{1}{c^2}\frac{\partial^2}{\partial t^2} - \nabla^2
\end{equation}


They are generated by accelerating mass distributions with a changing quadrupole moment.

\section*{Q3. Wavelength, frequency, wave speed, and radiation energy}

The \textbf{wavelength} $\lambda$ is the spatial distance between two consecutive points of the wave that have the same phase, such as crest to crest or trough to trough.

The \textbf{frequency} $\nu$, sometimes written $f$, is the number of oscillations per unit time.

The \textbf{period} $T$ is the time for one full oscillation:

\begin{equation}
T=\frac{1}{\nu}.
\end{equation}

The angular frequency is

\begin{equation}
\omega=2\pi\nu.
\end{equation}

The wave number is

\begin{equation}
k=\frac{2\pi}{\lambda}.
\end{equation}

For a sinusoidal traveling wave,

\begin{equation}
y(x,t)=A\cos(kx-\omega t+\phi),
\end{equation}

where $A$ is amplitude and $\phi$ is a phase constant.

The wave speed is

\begin{equation}
v=\lambda\nu.
\end{equation}

For electromagnetic radiation in vacuum,

\begin{equation}
c=\lambda\nu.
\end{equation}

The energy of a photon is related to frequency by Planck's relation:

\begin{equation}
E=h\nu.
\end{equation}

Since $\nu=c/\lambda$, the energy can also be written as

\begin{equation}
E=\frac{hc}{\lambda}.
\end{equation}

So higher-frequency radiation has higher energy, and shorter-wavelength radiation has higher energy.

The photon momentum is

\begin{equation}
p=\frac{E}{c}=\frac{h}{\lambda}.
\end{equation}

Therefore,

\begin{equation}
E=pc
\end{equation}

for massless radiation such as photons.

For material particles, the de Broglie relation is

\begin{equation}
\lambda=\frac{h}{p}.
\end{equation}

In relativistic form,

\begin{equation}
E^2=p^2c^2+m^2c^4.
\end{equation}

This reduces to $E=pc$ only when $m=0$.

\section*{Q4. Fundamental forces and their mediators}

Modern physics recognizes four fundamental interactions:

\begin{enumerate}
    \item Strong nuclear force
    \item Electromagnetic force
    \item Weak nuclear force
    \item Gravitational force
\end{enumerate}

They differ in strength, range, mediator particles, and the particles that experience them.

\subsection*{1. Strong interaction}

The \textbf{strong force} is described by quantum chromodynamics, or QCD. It acts on particles with \textbf{color charge}, especially quarks and gluons.

The mediator particles are \textbf{gluons}:

\[
g.
\]

The QCD gauge group is

\begin{equation}
SU(3)_C.
\end{equation}

The interaction term in the QCD Lagrangian has the form

\begin{equation}
\mathcal{L}_{\text{int}}=-g_s\bar{q}\gamma^\mu T^a qG^a_\mu,
\end{equation}

where $g_s$ is the strong coupling constant, $q$ is a quark field, $T^a$ are the generators of $SU(3)$, and $G^a_\mu$ is the gluon field.

Particles that participate in the strong interaction are

\[
\text{quarks and gluons.}
\]

Composite particles such as protons, neutrons, mesons, and nuclei experience the strong force because they are made of quarks and gluons.

The strong force is the strongest of the four at nuclear scales, but it is short-ranged due to confinement.

\subsection*{2. Electromagnetic interaction}

The electromagnetic force acts on particles with electric charge.

The mediator particle is the \textbf{photon}:

\[
\gamma.
\]

Its gauge group is

\begin{equation}
U(1)_{\text{EM}}.
\end{equation}

The electromagnetic interaction term is

\begin{equation}
\mathcal{L}_{\text{int}}=-e\bar{\psi}\gamma^\mu A_\mu\psi,
\end{equation}

where $e$ is the electric charge, $\psi$ is a charged fermion field, and $A_\mu$ is the electromagnetic four-potential.

Particles that participate are

\[
\text{charged leptons, quarks, } W^\pm \text{ bosons, and charged composite particles.}
\]

Examples include electrons, muons, taus, protons, charged pions, and atomic nuclei.

The electromagnetic force has infinite range because the photon is massless.

\subsection*{3. Weak interaction}

The weak force is responsible for processes such as beta decay, neutrino interactions, and flavor-changing reactions.

The mediator particles are

\[
W^+,\quad W^-,\quad Z^0.
\]

The electroweak gauge group before symmetry breaking is

\begin{equation}
SU(2)_L \times U(1)_Y.
\end{equation}

A charged-current weak interaction has the schematic form

\begin{equation}
\mathcal{L}_{\text{CC}}=-\frac{g}{\sqrt{2}}
\left(J^\mu_+W^+_\mu+J^\mu_-W^-_\mu\right),
\end{equation}

and the neutral-current interaction has the form

\begin{equation}
\mathcal{L}_{\text{NC}}=-\frac{g}{\cos\theta_W}J^\mu_Z Z_\mu.
\end{equation}

Particles that participate are

\[
\text{all quarks and leptons.}
\]

Neutrinos participate in the weak interaction but not in electromagnetic or strong interactions.

The weak force is short-ranged because $W^\pm$ and $Z^0$ are massive.

\subsection*{4. Gravitational interaction}

Gravity acts on all particles and fields with energy, momentum, and stress.

In general relativity, gravity is described not as an ordinary force but as curvature of spacetime. The source is the stress-energy tensor $T_{\mu\nu}$.

Einstein's field equations are

\begin{equation}
G_{\mu\nu}=\frac{8\pi G}{c^4}T_{\mu\nu}.
\end{equation}

In a quantum-field-theory description, the hypothetical mediator is the \textbf{graviton}:

\[
G \quad \text{or} \quad h_{\mu\nu}.
\]

The graviton is expected to be massless and spin-2, although it has not been directly detected as an individual quantum.

Particles that participate are

\[
\text{all particles with energy-momentum.}
\]

Gravity has infinite range, but it is by far the weakest interaction at the particle scale.

\subsection*{Order of strength}

At typical particle or nuclear scales, the approximate ordering is

\begin{equation}
\text{strong} > \text{electromagnetic} > \text{weak} > \text{gravity}.
\end{equation}

A rough relative comparison is

\begin{equation}
1 : 10^{-2} : 10^{-5} : 10^{-38}.
\end{equation}

This ordering is scale-dependent because coupling constants run with energy, but it is the usual low-energy/nuclear-scale ordering.

\begin{table}[h!]
\centering
\begin{tabular}{0.8\textwidth}
\toprule
\textbf{Force} & \textbf{Relative strength} & \textbf{Mediator} & \textbf{Range} & \textbf{Acts on} \\
\midrule
Strong & $1$ & Gluons & $\sim 10^{-15}\,\mathrm{m}$ & Quarks, gluons \\
Electromagnetic & $\sim 10^{-2}$ & Photon & Infinite & Electrically charged particles \\
Weak & $\sim 10^{-5}$ & $W^\pm, Z^0$ & $\sim 10^{-18}\,\mathrm{m}$ & Quarks and leptons \\
Gravity & $\sim 10^{-38}$ & Graviton, hypothetical & Infinite & All energy-momentum \\
\bottomrule
\end{tabular}
\caption{Approximate comparison of the four fundamental interactions.}
\label{tab:forces}
\end{table}

\section*{Q5. Fundamental particles in nature and their interactions}

In the Standard Model, fundamental particles are divided into:

\begin{enumerate}
    \item \textbf{Fermions}, which make up matter.
    \item \textbf{Bosons}, which mediate interactions.
    \item The \textbf{Higgs boson}, associated with electroweak symmetry breaking and particle masses.
\end{enumerate}

\subsection*{Fundamental fermions}

There are \textbf{12 fundamental fermions}, not counting antiparticles.

They consist of:

\begin{itemize}
    \item 6 quarks
    \item 6 leptons
\end{itemize}

Each has a corresponding antiparticle, so including antiparticles gives 24 fermionic particle species.

\subsection*{Quarks}

The six quarks are

\[
u,\ d,\ c,\ s,\ t,\ b.
\]

They are arranged in three generations.

\begin{table}[h!]
\centering
\begin{tabular}{lll}
\toprule
\textbf{Generation} & \textbf{Up-type quark} & \textbf{Down-type quark} \\
\midrule
1 & Up $u$ & Down $d$ \\
2 & Charm $c$ & Strange $s$ \\
3 & Top $t$ & Bottom $b$ \\
\bottomrule
\end{tabular}
\caption{The six quarks arranged by generation.}
\label{tab:quarks}
\end{table}

Quarks participate in strong, electromagnetic, weak, and gravitational interactions and they
carry color charge, electric charge, weak isospin/hypercharge, and energy-momentum.

\subsection*{Leptons}

The six leptons are

\[
e,\ \mu,\ \tau,\ \nu_e,\ \nu_\mu,\ \nu_\tau.
\]

They are arranged in three generations.

\begin{table}[h!]
\centering
\begin{tabular}{lll}
\toprule
\textbf{Generation} & \textbf{Charged lepton} & \textbf{Neutrino} \\
\midrule
1 & Electron $e^-$ & Electron neutrino $\nu_e$ \\
2 & Muon $\mu^-$ & Muon neutrino $\nu_\mu$ \\
3 & Tau $\tau^-$ & Tau neutrino $\nu_\tau$ \\
\bottomrule
\end{tabular}
\caption{The six leptons arranged by generation.}
\label{tab:leptons}
\end{table}

Charged leptons participate in electromagnetic, weak, and gravitational interactions.

Neutrinos participate in weak and gravitational interactions.

Neutrinos do not participate in the strong interaction and have no electric charge.

\subsection*{Fundamental bosons}

The gauge bosons are

\[
\gamma,\quad g,\quad W^+,\quad W^-,\quad Z^0.
\]

More specifically:

\begin{itemize}
    \item 1 photon
    \item 8 gluons
    \item 3 weak bosons: $W^+$, $W^-$, and $Z^0$
\end{itemize}

So the Standard Model has 12 gauge bosons if the eight gluon color states are counted separately.

The Higgs boson is

\[
H^0.
\]

Thus, in the Standard Model, the usual count is

\begin{equation}
6\text{ quarks}+6\text{ leptons}+12\text{ gauge bosons}+1\text{ Higgs}=25
\end{equation}

fundamental particles, not counting antiparticles separately.

If antiparticles are counted separately, the number increases.

\section*{Particles participating in each interaction}

\subsection*{Strong interaction}

Particles participating directly are

\[
q,\ \bar{q},\ g.
\]

That means all quarks, antiquarks, and gluons.

Leptons do not experience the strong force.

\subsection*{Electromagnetic interaction}

Particles participating directly are those with nonzero electric charge:

\[
e^-,\ \mu^-,\ \tau^-,
\]

\[
u,\ d,\ c,\ s,\ t,\ b,
\]

\[
W^+,\ W^-,
\]

and their antiparticles.

The photon mediates the electromagnetic interaction, but photons do not carry electric charge, so they do not directly couple to other photons at tree level in ordinary QED.

Neutrinos do not participate electromagnetically because

\begin{equation}
Q_{\nu}=0.
\end{equation}

\subsection*{Weak interaction}

All quarks and leptons participate in the weak interaction:

\[
u,\ d,\ c,\ s,\ t,\ b,
\]

\[
e,\ \mu,\ \tau,\ \nu_e,\ \nu_\mu,\ \nu_\tau.
\]

The weak interaction can change flavor through charged-current interactions. For example, beta decay is mediated by the $W^-$ boson:

\begin{equation}
d \rightarrow u + W^-,
\end{equation}

\begin{equation}
W^- \rightarrow e^-+\bar{\nu}_e.
\end{equation}

Together,

\begin{equation}
d \rightarrow u+e^-+\bar{\nu}_e.
\end{equation}

Inside a neutron,

\begin{equation}
n \rightarrow p+e^-+\bar{\nu}_e.
\end{equation}

Neutral-current weak interactions are mediated by the $Z^0$ boson and do not change electric charge.

\subsection*{Gravitational interaction}

All particles participate gravitationally because all have energy and momentum.

The gravitational coupling is to the stress-energy tensor:

\[
T_{\mu\nu}.
\]

Thus gravity acts on

\[
\text{quarks, leptons, gauge bosons, Higgs bosons, photons, neutrinos, and composite particles.}
\]

\section*{Summary table}

\begin{table}[h!]
\centering
\begin{tabular}{0.8\textwidth}
\toprule
\textbf{Interaction} & \textbf{Mediator} & \textbf{Participating fundamental particles} \\
\midrule
Strong & Gluons $g$ & Quarks, antiquarks, gluons \\
Electromagnetic & Photon $\gamma$ & Charged leptons, quarks, $W^\pm$, charged antiparticles \\
Weak & $W^\pm, Z^0$ & All quarks and leptons \\
Gravity & Graviton, hypothetical & All particles with energy-momentum \\
\bottomrule
\end{tabular}
\caption{Fundamental interactions and participating particles.}
\label{tab:interaction_summary}
\end{table}

The Standard Model contains

\begin{equation}
12\ \text{fermions} + 12\ \text{gauge bosons} + 1\ \text{Higgs boson} = 25
\end{equation}

fundamental particle types, excluding antiparticles as separate entries.

\end{document}